\documentclass[12pt]{article}
\usepackage[utf8]{inputenc}
\usepackage[T1]{fontenc}
\usepackage{hyperref}
\title{Taloki Language Resources for K-12 Curriculum Development}
\author{Compiled by Claude Sonnet 5 (claude-sonnet-5) for Andr\'as Kornai}
\date{}
\begin{document}
\maketitle

\section{Introduction}
Taloki (ISO 639-3 tlk) is an endangered Austronesian language of Buton Island, Southeast Sulawesi, Indonesia. It is spoken mainly around Maligano. It belongs to the Bungku-Tolaki group and is close to Kulisusu and Koroni. Speaker estimates are about 550 to 600. One study reports language shift among both older and younger people. No educational materials were found.

\section{Primary Educational Resources}
None found in this search.

\section{Linguistic and Reference Materials}
\begin{itemize}
\item Austronesian Basic Vocabulary Database entry for Taloki (word list): \url{https://abvd.eva.mpg.de/austronesian/language.php?id=908}
\item Verbix page on Taloki, with lexical similarity figures: \url{https://docs.verbix.com/Languages/Taloki}
\item Lexicostatistics article relating Taloki and Tolaki (Atlantis Press): \url{https://www.atlantis-press.com/article/125967751.pdf}
\end{itemize}
David Mead's 1999 monograph on the Bungku-Tolaki languages was named in the search summary, but I saw no URL for it.

\section{Religious and Community Resources}
\begin{itemize}
\item Joshua Project profile of the Taloki people: \url{https://joshuaproject.net/people_groups/15219}
\end{itemize}

\section{Digital and Audio Resources}
None found.

\section{Curriculum Development Context}
There is a word list and some comparative linguistics but no literacy materials. Any curriculum would start from documentation work with the remaining speakers, mostly older adults. Indonesian is the likely language of instruction.

\end{document}